% Fragmento de anexo. El documento receptor decide dónde ejecutar \appendix.
\section{Auditoría del concentrador y memoria de cálculo M8}
\label{anexo:Kt_M8}
\subsection{Recuperación del procedimiento de EP2}
Se revisó el documento original \texttt{Entrega\_2.pdf}: página PDF 12,
referencia [11], y páginas PDF 19--20, apartados A.4.5 y A.4.6.
La referencia general es Budynas y Nisbett, \emph{Diseño en ingeniería
mecánica de Shigley}, novena edición, 2012. El cuadro 18 usa
$K_t=1$, $1,5$ y $2,5$ como barrido paramétrico. No se identifica
una tabla de roscas, gráfico, interpolación, radio de raíz ni ecuación
geométrica que produzca estos valores\cite{ep2_auditoria_Kt}.

El código original \texttt{estructural/resistencia\_fatiga.py}, bloque
\texttt{ENTALLAS}, confirma: escenario moderado $K_t=1,5$, $q=0,70$,
declarado no calibrado; escenario alto $K_t=2,5$, $q=0,90$, declarado
no validado. Las ecuaciones 35--38 de EP2 corresponden a Marin,
resistencia inicial, sensibilidad a la entalla y Goodman:
\begin{align}
 S_e&=k_a k_b k_c k_d k_e k_{\mathrm{misc}} S'_e,\\
 K_f&=1+q(K_t-1),\\
 \frac{1}{n_f}&=\frac{K_f\sigma_a}{S_e}
 +\frac{K_t\max(\sigma_m,0)}{S_{ut}}.
\end{align}
Estas expresiones no determinan $K_t$. Ni los valores de $q$ ni los
factores de Marin de EP2 se transfieren automáticamente al producto M8.
La memoria anterior se conserva como antecedente; esta revisión limita
su interpretación, sin modificar sus tablas ni sus archivos.

\subsection{Fuentes consultadas y límites de aplicación}
\begin{description}
\item[ISO 68-1:2023.]
Se consultó la vista previa, páginas impresas 1--2: distingue el perfil
básico del perfil de diseño y sus diámetros $d_1$ y $d_3$
\cite{iso68_2023_vista_Kt}. La designación M8 no define por sí sola el
perfil fabricado y su tolerancia. La vista previa no incluye toda la
geometría de diseño; no se afirma lectura completa de la norma.
\item[ISO 965-1.]
La edición 2013 fue reemplazada por ISO 965-1:2026, según el catálogo ISO
\cite{iso965_2026_estado_Kt}. La vista previa 2026 identifica el apartado
12 de contornos de raíz y el anexo B de diámetros menores, pero no permite
verificar aquí sus requisitos completos. Como antecedente, el contenido
indexado del apartado 11 de la edición 2013 indica un radio mínimo
$0,125P$ para clases 8.8 o superiores y una raíz de referencia
$R\simeq0,14434P$\cite{iso965_2013_antecedente_Kt}.
Para $P=1,25$ mm, serían 0,15625 mm y aproximadamente 0,18042 mm.
El primero no se presenta como requisito 2026 comprobado y ninguno
se asigna como radio real. Es indispensable confirmar la edición
contractual y el perfil del producto.
\item[Pedersen, 2013.]
El resumen del repositorio del autor estudia roscas métricas ISO,
geometría de raíz y contacto mediante un modelo axisimétrico de unión
roscada\cite{pedersen_2013_roscas}. Sustenta la necesidad de definir
raíz/contacto; no proporciona un $K_t$ directamente transferible a la
flexión global de esta varilla.
\item[Ohara, 1980.]
La ficha NASA describe concentración mediante transferencia de carga
entre filetes y contacto de flancos\cite{ohara_1980_roscas}.
No se usa una ordenada de ese trabajo como concentrador de flexión libre.
\item[Brutti y colaboradores, 2025.]
El apéndice A.1 trata esfuerzos locales del primer filete engranado;
su factor de ``flexión'' multiplica la flexión del diente, con su propio
módulo resistente\cite{brutti_2025_roscas}. No equivale a multiplicar
$32M/(\pi d_r^3)$ por ese factor. Se descarta esa sustitución.
\item[Peterson/Pilkey.]
El catálogo editorial de la cuarta edición identifica la sección 5.19
sobre piezas roscadas\cite{peterson_2020_catalogo_Kt}. No se tuvo acceso
verificable a sus gráficos completos; no se declara lectura ni selección
de una curva o valor numérico de esa obra. Para emplearla será necesario
disponer del gráfico, verificar carga, geometría y denominador nominal.
\end{description}

\subsection{Sección nominal y cambio de base del concentrador}
Con $d_r=6,466$ mm y $D=8$ mm se define, para un mismo momento global:
\begin{align}
 I_r&=\frac{\pi d_r^4}{64}, & Z_r&=\frac{\pi d_r^3}{32},\\
 \sigma_{\mathrm{nom},r}&=\frac{M}{Z_r}, &
 \sigma_{\mathrm{nom},D}&=\sigma_{\mathrm{nom},r}\left(\frac{d_r}{D}\right)^3.
\end{align}
La igualdad de tensiones locales exige
\begin{equation}
 K_{t,D}=K_{t,r}\left(\frac{D}{d_r}\right)^3,
 \qquad
 K_{t,r}=K_{t,D}\left(\frac{d_r}{D}\right)^3.
 \label{eq:conversion_Kt_bases}
\end{equation}
Esta conversión algebraica es válida para las dos referencias circulares
indicadas y el mismo momento. No convierte un concentrador de carga axial
o flexión de diente en uno de flexión global. Si una fuente ya define
su nominal con el diámetro menor, \textbf{no se aplica otra corrección
por diámetro}. Si usa otro módulo resistente, se utiliza ese módulo,
no se fuerza la conversión anterior.

El archivo \texttt{resumen.json} conserva los valores sin redondear de
$I$, momento y cambio de base mediante las entradas y ecuaciones de esta
memoria; \texttt{sensibilidad\_Kt\_sin\_seleccion.csv} registra ambos
denominadores. La tensión nominal sobre el núcleo no es una solución
local de la hélice: el STEP tiene una representación cilíndrica y no
resuelve el radio de raíz. El cilindro equivalente de diámetro de núcleo
es una idealización de rigidez; no se demuestra que sus tensiones sean
una cota conservadora de todas las tensiones locales.
El área resistente convencional de tracción de una rosca no sustituye
al módulo resistente de flexión ni a la inercia de la sección; tampoco
se adopta un diámetro equivalente derivado de esa área como una nueva
sección flexional sin demostrar su correspondencia geométrica.

En concreto, $(D/d_r)^3=1,893925$ y
$\sigma_{\mathrm{nom},D}=166,5123$ MPa. Un valor publicado sobre diámetro
mayor debe dividirse por 1,893925 antes de multiplicar la nominal de
315,361913 MPa, siempre que describa la misma flexión y geometría.
El umbral $K_{t,r}=2,029414$ equivale a $K_{t,D}=3,843559$; son dos
definiciones del mismo umbral, no dos predicciones diferentes.

La tensión máxima heredada es 315,361913 MPa en el estado 541, a
$x=0,128117780089$ m desde el extremo de referencia. En ese mismo estado:
$F=917,9576$ N, torque de accionamiento $2,63459$ N\,m y
flecha máxima $0,1568706$ mm. El momento resultante de 8,369822 N\,m
reproduce $\sigma=32M/(\pi d_r^3)$. El estado crítico está en una
fijación ideal: la raíz libre, la primera espira de una tuerca/inserto y
la salida de rosca pueden tener concentradores y cargas distintos.

\subsection{Sensibilidad estática, sin selección de factor}
Se conserva el límite mínimo candidato de 640 MPa de la clase 8.8,
sujeto a certificado de producto\cite{boellhoff_manual_8100_2023}.
\begin{center}
\begin{tabular}{rrrr}
\toprule
$K_{t,r}$ hipotético & $K_{t,r}\sigma_{\mathrm{nom},r}$ [MPa]
 & $n_{y,\mathrm{flex}}$ & Procedencia\\
\midrule
1,000 & 315,36 & 2,029 & Referencia nominal\\
1,500 & 473,04 & 1,353 & Sensibilidad EP2\\
2,000 & 630,72 & 1,015 & Diagnóstico matemático\\
2,029414 & 640,00 & 1,000 & Umbral calculado\\
2,500 & 788,40 & 0,812 & Sensibilidad EP2\\
\bottomrule
\end{tabular}
\end{center}
\textbf{Estos puntos no constituyen un intervalo bibliográfico validado
del $K_t$ de una M8.} Las fuentes sustentan la dependencia geométrica y
la necesidad de distinguir los modelos; EP2 solo aporta antecedentes
paramétricos. No se encontró una curva verificable aplicable, con el
perfil y las condiciones de montaje disponibles, que permita justificar
un intervalo numérico físico más estrecho. No se inventa ese intervalo.
Los radios geométricos de referencia tampoco permiten interpolar $K_t$
sin una curva o solución válida de flexión global.

El criterio $n_y\geq1$ exige $K_{t,r}\leq2,029414$ solamente para la
aproximación elástica de flexión considerada. Para objetivos hipotéticos
$n_y=1,5$ o $2$, los límites serían $1,352943$ y $1,014707$,
respectivamente; no son requisitos establecidos del proyecto.
No se suman el cortante máximo global y la flexión máxima si pertenecen
a puntos diferentes. Tampoco se transforma el torque vertical en torsión
del eje transversal. Una comprobación final requiere el tensor local
correlacionado, incluyendo las acciones reales de la fijación; puede
requerir von Mises y un tratamiento elastoplástico si hay fluencia local.

\subsection{Datos para identificar $K_t$ y cerrar fatiga}
Se necesita definir y verificar: diámetro de raíz mínimo y real, perfil
y radio de curvatura de raíz, paso, tolerancias y recubrimiento; fase
helicoidal y salida de rosca respecto de ambos empotramientos; tuercas o
insertos, longitud engranada, primer filete activo y precarga; transiciones,
asientos y forma de transmitir momento a la carcasa. Con ello se podrá
seleccionar una solución bibliográfica compatible o preparar un submodelo
elástico de la geometría local con convergencia de malla en el radio,
cargas correlacionadas y comparación contra la misma tensión nominal.
No se ejecutó ese submodelo en esta revisión.

$K_t$ es geométrico y elástico para una geometría, carga y nominal
definidos. $K_f$ expresa la reducción efectiva de resistencia a fatiga;
la relación $K_f=1+q(K_t-1)$ requiere una sensibilidad $q$ justificada.
El radio solo no determina $q$: faltan condición metalúrgica, fabricación
y superficie. La clase 8.8 no garantiza por sí sola $S_e$, $q$ ni una
curva S--N del producto. Si una resistencia cíclica de catálogo ya incluye
la rosca, no se multiplica nuevamente por un $K_f$ por el mismo efecto.

Se conservan los historiales firmados y los ciclos rainflow anteriores;
no se recalcularon. Para cerrar la vida faltan producto certificado,
rosca laminada o cortada y secuencia de tratamientos, estado superficial,
recubrimiento/tensiones residuales, $S_e$ o curva S--N aplicable,
confiabilidad, ambiente y vida/espectro de carga. Goodman permanece
condicional; no se asigna un factor de seguridad a fatiga ni se adopta
automáticamente $S'_e=0,5S_{ut}$ como propiedad garantizada.

\subsection{Trazabilidad y alcance de la revisión}
Las entradas son los CSV y configuraciones originales enumerados en
\texttt{historicos\_sha256.json}. El script local
\texttt{reproducir\_memoria.py} solo lee el historial; no importa ni llama
al modelo 3C. Se comprueban los SHA256 antes y después del postproceso.
Las tablas completas de esta revisión se almacenan separadas de las
memorias anteriores. Esta revisión no certifica la clase comprada,
los apoyos ABS, el montaje 608, la holgura admisible ni la resistencia
de la unión real. M8 continúa como alternativa de diseño preferente.
